\documentclass[12pt,a4paper]{article}

% ============================================
% PACKAGES
% ============================================
\usepackage[utf8]{inputenc}
\usepackage[T1]{fontenc}
\usepackage[french]{babel}
\usepackage{geometry}
\usepackage{graphicx}
\usepackage{hyperref}
\usepackage{booktabs}
\usepackage{enumitem}
\usepackage{titlesec}
\usepackage{fancyhdr}
\usepackage{xcolor}
\usepackage{tocloft}
\usepackage{parskip}

% ============================================
% CONFIGURATION
% ============================================
\geometry{
    left=2.5cm,
    right=2.5cm,
    top=2.5cm,
    bottom=2.5cm
}

\hypersetup{
    colorlinks=true,
    linkcolor=blue!60!black,
    urlcolor=blue!60!black,
    pdftitle={Cahier des Charges - Extraction automatique de réseaux de personnages},
    pdfauthor={Lotfi ABDALLAH, Wael MANSOURA}
}

% En-têtes et pieds de page
\pagestyle{fancy}
\fancyhf{}
\fancyhead[L]{\small Cahier des Charges}
\fancyhead[R]{\small Réseaux de Personnages}
\fancyfoot[C]{\thepage}
\renewcommand{\headrulewidth}{0.4pt}
\renewcommand{\footrulewidth}{0.4pt}

% Formatage des sections
\titleformat{\section}
    {\normalfont\Large\bfseries\color{blue!60!black}}
    {\thesection.}{0.5em}{}
\titleformat{\subsection}
    {\normalfont\large\bfseries}
    {\thesubsection.}{0.5em}{}

% ============================================
% DOCUMENT
% ============================================
\begin{document}

% ============================================
% PAGE DE TITRE
% ============================================
\begin{titlepage}
    \centering
    \vspace*{2cm}
    
    {\scshape\LARGE Université d'Avignon \par}
    \vspace{0.5cm}
    {\scshape\Large Master 1 -- AMS \par}
    \vspace{2cm}
    
    \rule{\textwidth}{1.5pt}\par
    \vspace{0.8cm}
    {\Huge\bfseries Cahier des Charges\par}
    \vspace{0.5cm}
    {\LARGE Extraction automatique de réseaux de personnages\par}
    \vspace{0.8cm}
    \rule{\textwidth}{1.5pt}\par
    
    \vspace{2cm}
    
    {\Large\itshape Lotfi ABDALLAH \\ Wael MANSOURA\par}
    
    \vfill
    
    {\large Année universitaire 2024--2025\par}
\end{titlepage}

% ============================================
% TABLE DES MATIÈRES
% ============================================
\tableofcontents
\newpage

% ============================================
% SECTION 1 : CONTEXTE ET OBJECTIFS
% ============================================
\section{Contexte et Objectifs du Projet}

\subsection{Contexte}

L'extraction automatique de réseaux de personnages à partir de textes narratifs est une tâche complexe qui combine des techniques de traitement du langage naturel (NLP) et la modélisation graphique.

Ce projet vise à développer un système capable~:
\begin{itemize}[noitemsep]
    \item d'identifier les personnages dans un texte,
    \item de détecter leurs interactions,
    \item de construire un réseau représentatif des relations entre ces personnages.
\end{itemize}

\subsection{Objectifs}

Les objectifs principaux du projet sont~:
\begin{itemize}
    \item Développer un algorithme d'extraction de personnages à partir de textes narratifs.
    \item Construire des listes~:
    \begin{itemize}
        \item \textbf{L} -- toutes les entités nommées candidates,
        \item \textbf{LP} -- les personnages (personnes),
        \item \textbf{LL} -- les lieux.
    \end{itemize}
    \item Détecter les cooccurrences de personnages en utilisant la proximité textuelle.
    \item Construire un réseau de personnages basé sur les interactions détectées.
\end{itemize}

% ============================================
% SECTION 2 : PÉRIMÈTRE DU PROJET
% ============================================
\section{Périmètre du Projet}

\subsection{Fonctionnalités attendues}

\begin{enumerate}
    \item Extraction des entités nommées.
    \item Filtrage des personnes (liste \textbf{LP}).
    \item Filtrage des lieux (liste \textbf{LL}).
    \item Regroupement d'alias pour chaque personnage.
    \item Génération du graphe avec \texttt{NetworkX} pour chaque chapitre.
    \item Détection des interactions entre les personnages.
    \item Construction et visualisation du réseau de personnages.
    \item Export CSV des résultats pour plusieurs chapitres.
\end{enumerate}

\subsection{Technologies utilisées}

\begin{table}[h!]
\centering
\begin{tabular}{@{}ll@{}}
\toprule
\textbf{Catégorie} & \textbf{Outils / Technologies} \\
\midrule
Langage de programmation & Python \\
Bibliothèques NLP & spaCy, Stanza, Flair \\
Visualisation & NetworkX, Matplotlib \\
Filtrage & \texttt{antidict.txt} (entités non pertinentes) \\
\bottomrule
\end{tabular}
\caption{Technologies et outils utilisés}
\end{table}

% ============================================
% SECTION 3 : PUBLIC CIBLE
% ============================================
\section{Public Cible}

\subsection{Profil des utilisateurs}

\begin{itemize}
    \item Étudiants en informatique.
    \item Chercheurs en linguistique computationnelle.
    \item Développeurs en traitement du langage naturel.
\end{itemize}

\subsection{Besoins des utilisateurs}

\begin{itemize}
    \item Outils efficaces pour l'extraction de réseaux de personnages.
    \item Visualisation claire des réseaux de personnages.
\end{itemize}

% ============================================
% SECTION 4 : SCÉNARIOS DE TESTS UTILISATEURS
% ============================================
\section{Scénarios de Tests Utilisateurs}

\subsection{Objectifs des tests}

\begin{itemize}
    \item Valider que les outils d'extraction identifient correctement les personnages importants.
    \item Vérifier la pertinence du regroupement d'alias.
    \item Évaluer la cohérence des co-occurrences détectées.
\end{itemize}

\subsection{Méthodologie}

\begin{enumerate}
    \item Sélectionner un ensemble de textes narratifs variés.
    \item Appliquer les outils d'extraction développés.
    \item Comparer les résultats obtenus avec des annotations manuelles.
\end{enumerate}

% ============================================
% SECTION 5 : CRITÈRES DE RÉUSSITE
% ============================================
\section{Critères de Réussite et Évaluations}

\subsection{Indicateurs de performance (KPI)}

\begin{itemize}
    \item Taux de précision de l'extraction des personnages.
    \item Score de précision du graphe de personnages par chapitre sur le classement Kaggle.
    \item Temps de génération du graphe pour chaque chapitre.
\end{itemize}

\subsection{Feedback des utilisateurs / professeurs}

\begin{itemize}
    \item Utilisation de plusieurs modèles NLP pour extraire les personnages.
    \item Méthode de \textbf{vote} pour identifier les personnages corrects, plutôt que d'utiliser l'intersection ou l'union simple des résultats.
\end{itemize}

% ============================================
% SECTION 6 : CONTRAINTES ET RISQUES
% ============================================
\section{Contraintes et Risques}

\subsection{Contraintes techniques}

\begin{itemize}
    \item Limites des modèles NLP pour la reconnaissance des entités nommées.
    \item Gestion des ambiguïtés dans les noms des personnages.
    \item Performance et scalabilité pour traiter de grands textes narratifs.
    \item Fusion des résultats de plusieurs modèles NLP pour améliorer la précision.
\end{itemize}

\subsection{Risques potentiels}

\begin{itemize}
    \item Mauvaise détection d'entités à cause du style littéraire.
    \item Difficulté à atteindre un taux de précision élevé.
    \item Problèmes de performance avec de grands volumes de texte.
\end{itemize}

% ============================================
% SECTION 7 : PLANNING PRÉVISIONNEL
% ============================================
\section{Planning Prévisionnel}

\subsection{Phases du projet}

\begin{table}[h!]
\centering
\begin{tabular}{@{}lc@{}}
\toprule
\textbf{Phase} & \textbf{Durée} \\
\midrule
Analyse et traitement du corpus & 2 semaines \\
Développement extraction NER \& listes L/LP/LL & 3 semaines \\
Détection des co-occurrences & 2 semaines \\
Génération des graphes & 2 semaines \\
Ajustements finaux + rapport & 1 semaine \\
\midrule
\textbf{Total} & \textbf{10 semaines} \\
\bottomrule
\end{tabular}
\caption{Planning prévisionnel du projet}
\end{table}

\subsection{Livrables}

\begin{itemize}
    \item Listes L, LP, LL.
    \item Graphes de personnages par chapitre.
    \item Fichiers CSV des résultats.
    \item Cahier des charges finalisé.
\end{itemize}

% ============================================
% SECTION 8 : ÉQUIPE DU PROJET
% ============================================
\section{Équipe du Projet}

\begin{table}[h!]
\centering
\begin{tabular}{@{}ll@{}}
\toprule
\textbf{Membre} & \textbf{Responsabilités} \\
\midrule
Lotfi ABDALLAH & Développement NLP, extraction NER, listes L/LP/LL \\
Wael MANSOURA & Détection co-occurrences, génération graphes \\
\bottomrule
\end{tabular}
\caption{Répartition des responsabilités}
\end{table}

% ============================================
% SECTION 9 : ANNEXES
% ============================================
\section{Annexes}

\subsection{Architecture du système}

Le pipeline d'extraction de réseaux de personnages suit l'architecture suivante~:

\begin{center}
\fbox{%
\begin{minipage}{0.9\textwidth}
\centering
\vspace{0.5em}
\textbf{Pipeline de traitement}
\vspace{0.5em}

\texttt{Texte brut (.txt)}

$\downarrow$

\texttt{NER multi-modèles (spaCy + Stanza)}

$\downarrow$

\texttt{Liste L (toutes entités)}

$\downarrow$

\texttt{Filtrage + antidict.txt}

$\downarrow$

\texttt{Liste LP (personnes) \quad Liste LL (lieux)}

$\downarrow$

\texttt{Regroupement d'alias}

$\downarrow$

\texttt{Détection co-occurrences (fenêtre glissante)}

$\downarrow$

\texttt{Génération graphe NetworkX}

$\downarrow$

\texttt{Export (GraphML, CSV, visualisation)}
\vspace{0.5em}
\end{minipage}
}
\end{center}

\vspace{1em}

Le notebook \texttt{main.ipynb} orchestre l'ensemble du pipeline et permet une exécution interactive étape par étape.

\subsection{Description des scripts Python}

\begin{table}[h!]
\centering
\small
\begin{tabular}{@{}p{4.5cm}p{10cm}@{}}
\toprule
\textbf{Fichier} & \textbf{Rôle} \\
\midrule
\texttt{main.ipynb} & \textbf{Notebook principal} -- orchestre le pipeline complet : chargement des données, extraction NER, filtrage, alias, co-occurrences, génération et visualisation des graphes. Point d'entrée interactif du projet. \\
\addlinespace
\texttt{nlp\_main.py} & Script d'orchestration en ligne de commande. Appelle les modules dans l'ordre pour produire les listes L, LP, LL. \\
\addlinespace
\texttt{nlp\_extract\_characters.py} & Extraction des entités nommées avec spaCy. Fonctions \texttt{extract\_entities()}, \texttt{count\_entities()}, \texttt{filter\_persons()}, \texttt{filter\_locations()}. \\
\addlinespace
\texttt{nlp\_multi\_ner.py} & NER multi-modèles (spaCy + Stanza). Implémente la méthode de \textbf{vote} via \texttt{ensemble\_entities()} pour fusionner les résultats des différents modèles. \\
\addlinespace
\texttt{nlp\_aliases.py} & Regroupement d'alias : normalisation des noms, détection des variantes (ex: ``Hari'' $\rightarrow$ ``Hari Seldon''), fusion des comptages. \\
\addlinespace
\texttt{nlp\_cooccurrence.py} & Détection des co-occurrences entre personnages dans une fenêtre de mots configurable (\texttt{distance\_max}). \\
\addlinespace
\texttt{nlp\_graph.py} & Génération du graphe NetworkX, sauvegarde en GraphML, visualisation statique avec Matplotlib. \\
\addlinespace
\texttt{nlp\_visualize\_web.py} & Visualisation interactive du graphe en HTML avec PyVis. \\
\addlinespace
\texttt{nlp\_create\_submission.py} & Traitement par chapitre et génération du fichier \texttt{submission.csv} pour Kaggle. \\
\addlinespace
\texttt{nlp\_utils.py} & Fonctions utilitaires : lecture de fichiers, chargement de l'anti-dictionnaire. \\
\bottomrule
\end{tabular}
\caption{Description des modules Python du projet}
\end{table}

\end{document}
